\noindent
\textcolor{stewartblue}{\textsf{\textbf{VÍ DỤ 6}}}\quad Tính $\displaystyle \lim_{x \to 2^+} \arctan\left(\frac{1}{x - 2}\right)$.

\vspace{0.8ex}
\noindent
\textcolor{stewartblue}{\textsf{\textbf{LỜI GIẢI}}}\quad Nếu đặt $t = 1/(x - 2)$, ta biết rằng $t \to \infty$ khi $x \to 2^+$. Do đó, theo phương trình thứ hai trong (4), ta có
\[
\lim_{x \to 2^+} \arctan\left(\frac{1}{x - 2}\right) = \lim_{t \to \infty} \arctan t = \frac{\pi}{2} \tag*{\textcolor{stewartblue}{\rule{1.15ex}{1.15ex}}}
\]

\noindent
Đồ thị của hàm số mũ tự nhiên $y = e^x$ có đường thẳng $y = 0$ (trục $x$) là một tiệm cận ngang. (Điều này cũng đúng với mọi hàm số mũ có cơ số $b > 1$.) Thật vậy, từ đồ thị trong Hình 10 và bảng giá trị tương ứng, ta thấy rằng

\begin{center}
\makebox[\linewidth]{%
  \makebox[0pt][l]{\tikz{\node[draw=stewartred, semithick, inner sep=2.5pt, font=\bfseries\sffamily\color{stewartred}] {6};}}%
  \hfill
  \tikz{\node[draw=stewartred, semithick, inner sep=6pt, minimum width=2.8cm] {$\displaystyle \lim_{x \to -\infty} e^x = 0$};}%
  \hfill
  \phantom{\tikz{\node[draw=stewartred, semithick, inner sep=2.5pt] {6};}}%
}
\end{center}

\noindent
Lưu ý rằng các giá trị của $e^x$ tiến dần về $0$ rất nhanh.

\begin{center}
\marginnote{\textsf{\textbf{HÌNH 10}}}[2.8cm]%
\begin{minipage}[c]{0.52\linewidth}
  \centering
  \begin{tikzpicture}[scale=0.95, >=stealth]
    % Trục tọa độ
    \draw[->] (-3.8, 0) -- (2.2, 0) node[below=1pt] {$x$};
    \draw[->] (0, -0.4) -- (0, 3.8) node[left=1pt] {$y$};
    % Vạch chia
    \draw (1, 2pt) -- (1, -2pt) node[below=2pt, font=\small] {$1$};
    \draw (2pt, 1) -- (-2pt, 1) node[left=2pt, font=\small] {$1$};
    \node[below left=1pt, font=\small] at (0, 0) {$0$};
    % Đồ thị
    \draw[stewartblue, very thick, domain=-3.6:1.3, samples=100] plot (\x, {exp(\x)});
    % Nhãn
    \node[right, font=\small] at (0.6, 2.3) {$y = e^x$};
  \end{tikzpicture}
\end{minipage}\hfill
\begin{minipage}[c]{0.45\linewidth}
  \centering
  \renewcommand{\arraystretch}{1.15}
  \begin{tabular}{|>{\centering\arraybackslash}m{1.4cm}|>{\centering\arraybackslash}m{2.3cm}|}
    \hline
    \rowcolor{cyan!8}
    $x$ & $e^x$ \\
    \hline
    \phantom{$-$}0 & 1.00000 \\
    $-1$ & 0.36788 \\
    $-2$ & 0.13534 \\
    $-3$ & 0.04979 \\
    $-5$ & 0.00674 \\
    $-8$ & 0.00034 \\
    $-10$ & 0.00005 \\
    \hline
  \end{tabular}
\end{minipage}
\end{center}

\vspace{0.8ex}
\marginnote{%
  \parbox{44mm}{%
    \raggedright\textsf{\fontsize{7.5pt}{9.5pt}\selectfont
    \tikz[baseline=-0.6ex]{\node[fill=stewartpurple, text=white, font=\bfseries\sffamily\fontsize{7pt}{7pt}\selectfont, inner sep=1.8pt, rounded corners=0.5pt] {PS};}\;%
    Chiến lược giải quyết vấn đề cho các Ví dụ 6 và 7 là \textit{đưa vào một yếu tố phụ} (xem Các nguyên lý giải quyết vấn đề sau Chương 1). Ở đây, yếu tố phụ trợ chính là biến mới $t$.%
    }%
  }%
}[0pt]%
\noindent
\textcolor{stewartblue}{\textsf{\textbf{VÍ DỤ 7}}}\quad Tính $\displaystyle \lim_{x \to 0^-} e^{1/x}$.

\vspace{0.8ex}
\noindent
\textcolor{stewartblue}{\textsf{\textbf{LỜI GIẢI}}}\quad Nếu đặt $t = 1/x$, ta biết rằng $t \to -\infty$ khi $x \to 0^-$. Do đó, theo (6),
\[
\lim_{x \to 0^-} e^{1/x} = \lim_{t \to -\infty} e^t = 0
\]
\noindent
(Xem Bài tập 81.)\hfill\textcolor{stewartblue}{\rule{1.15ex}{1.15ex}}

\vspace{1.2ex}
\noindent
\textcolor{stewartblue}{\textsf{\textbf{VÍ DỤ 8}}}\quad Tính $\displaystyle \lim_{x \to \infty} \sin x$.

\vspace{0.8ex}
\noindent
\textcolor{stewartblue}{\textsf{\textbf{LỜI GIẢI}}}\quad Khi $x$ tăng, các giá trị của $\sin x$ dao động liên tục vô hạn lần giữa $1$ và $-1$, và do đó chúng không tiến đến bất kỳ một số xác định nào. Như vậy $\lim_{x \to \infty} \sin x$ không tồn tại.\hfill\textcolor{stewartblue}{\rule{1.15ex}{1.15ex}}

\vspace{2ex}
\noindent
\textcolor{stewartblue}{\rule{1.2ex}{1.2ex}}\hspace{0.5em}\textsf{\textbf{\large Giới hạn vô cực tại vô cực}}

\vspace{1ex}
\noindent
Ký hiệu
\[
\lim_{x \to \infty} f(x) = \infty
\]
được dùng để chỉ rằng các giá trị của $f(x)$ trở nên rất lớn khi $x$ trở nên rất lớn. Các ý nghĩa tương tự được quy ước cho các ký hiệu sau:
\[
\lim_{x \to -\infty} f(x) = \infty \qquad\qquad \lim_{x \to \infty} f(x) = -\infty \qquad\qquad \lim_{x \to -\infty} f(x) = -\infty
\]

\vfill
\begin{center}
\fontsize{5.5pt}{7pt}\selectfont\color{gray}
Copyright 2021 Cengage Learning. All Rights Reserved. May not be copied, scanned, or duplicated, in whole or in part. Due to electronic rights, some third party content may be suppressed from the eBook and/or eChapter(s).\\
Editorial review has deemed that any suppressed content does not materially affect the overall learning experience. Cengage Learning reserves the right to remove additional content at any time if subsequent rights restrictions require it.
\end{center}
